%!TEX root = ../template.tex
%%%%%%%%%%%%%%%%%%%%%%%%%%%%%%%%%%%%%%%%%%%%%%%%%%%%%%%%%%%%%%%%%%%%
%% abstract-en.tex
%% NOVA thesis document file
%%
%% Abstract in English
%%%%%%%%%%%%%%%%%%%%%%%%%%%%%%%%%%%%%%%%%%%%%%%%%%%%%%%%%%%%%%%%%%%%

\typeout{NT FILE abstract-en.tex}%

The Industry 4.0 paradigm relies on \glspl{CPS} and \glspl{CPPS} to achieve agile and resilient production. Within this landscape, \glspl{DT} provide virtual representations capable of mirroring physical assets through real-time bidirectional telemetry. However, traditional development relies on proprietary simulation software, fragmented 3D \gls{CAD} files lacking kinematics, and tedious manual matching of \gls{IO} addresses, making reconfigurations slow and error-prone.

To overcome these barriers, this dissertation presents an open, modular, model-driven framework for automated procedural generation and live synchronization of industrial \glspl{DT}. Developed under the \gls{DSR} methodology, the architecture spans four decoupled tiers: low-polygon asset digitization in Blender (Layer~0), no-code semantic authoring with 3D raycasting and \gls{AASX} packaging (Layer~1), containerized storage and discovery via Eclipse BaSyx~v2 microservices (Layer~2), and procedural scene assembly, physics simulation, dual-mode execution, and drift correction in Unity 3D (Layer~3).

The framework was empirically validated on a modular \gls{FMS} bench driven by an Industrial Shields \gls{PLC} running FreeRTOS, as well as through parametric stress testing across 800 simulation runs scaling from 5 to 1000 conveyor modules. Results demonstrate cold-boot instantiation of the physical pilot bench in under 2.8\,s and steady-state rendering rates exceeding 480\,\gls{FPS}. Scalability testing confirms near-linear scaling up to 250 modules and stable real-time execution up to 500 modules (averaging over 90\,\gls{FPS}), establishing model-driven \gls{AAS} standards as an effective foundation for next-generation \glspl{DT}.

\keywords{
  Industry 4.0 \and
  Cyber-Physical Systems \and
  Procedural Simulation \and
  Virtual Commissioning \and
  Digital Twin \and
  Asset Administration Shell
}